# Profil J3 — Prise de notes brute (Consultant senior, 3 ans)

**Code :** J3  
**Fonction :** Consultant senior  
**Ancienneté :** 3 ans  
**Contexte :** Cybersécurité, institution européenne. Beaucoup de documents Word, peu de détails sur les missions.

---

## T0 — Cadrage

**Q : Parcours, comment arrivé ici ?**
- 3 ans. Consultant senior.
- Institution européenne, cybersécurité. Très exigeant. Confidentialité, rigueur.
- Pas de détails sur les missions.
- Word. Livrables très normés. Procédure à respecter.

---

## T1 — Usage réel

**Q : Dernière fois utilisé outil génératif ?**
- Semaine dernière. Synthèse. 50 pages de notes techniques -> résumé pour comité de direction.
- Copilot. Filié doc. Demandé structure avec enjeux et reco.
- Proposé trame. 1h à retravailler, recouper avec notes. Vérifier déformations.
- Jamais confiance premier jet. Assistant, pas rédacteur.

**Q : Sur quelles tâches au quotidien ?**
- Trois choses :
  1. Reformulation technique -> accessible. Clients pas toujours niveau tech.
  2. Relecture. Incohérences, formulations ambigües.
  3. Premiers jets sur sujets moins maîtrisés.
- Jamais : architecture sécurité, vulnérabilités, données classifiées. Main humaine.

**Q : Comment appris ?**
- Tâtonné. Puis compris besoin structure. Bloc-notes avec prompts qui marchent pour tâches récurrentes.
- Synthèse, reformulation client, détection incohérences.
- Formation certifiante Microsoft depuis quelques mois. Pas juste apprendre, papier à valoriser. Demande clients forte.

**Q : Temps gagné mis sur quoi ?**
- Synthèse : gagné 2h sur structuration. Utilisées relecture croisée sources.
- IA fait erreurs subtiles. Formulations justes mais sens trahi.
- Préfère vérifier que produire. Qualité > vitesse.

**Q : Situation où plus d'effort que de le faire soi-même ?**
- Plusieurs fois. Plus marquante : reformulation technique. IA inversé 2 concepts clés. Sens faussé.
- 45min à reprendre. Aurais fait en 20min. Catastrophe si parti.
- Depuis, encore plus vigilant.

**Q : Vous en parlez entre vous ?**
- Oui, mais pas dans grandes largeurs. Retours d'expérience, mises en garde.
- Alerté juniors : pas confiance aveugle. Vu livrables avec coquilles grossières. Pas relu.
- Règle : "si tu peux pas justifier, mets pas". J'essaie de transmettre.

---

## T2 — Apprentissage

**Q : Comment on apprend le métier aujourd'hui ?**
- Expérience et erreurs. Mes premières missions, j'ai fait des erreurs.
- IA change donne : version propre -> illusion de maîtrise sans avoir creusé.
- Bon consultant = sait ce qu'il met dans livrables, pas formatage prompt.
- Formation continue me conforte : IA outil, pas compétence.

**Q : Par quelles tâches on apprend vraiment ?**
- Seul face au problème. Analyse, justifier reco, expliquer vulnérabilité.
- IA connaît pas contexte client, contraintes, sensibilités.
- Répétitives = pas formatrices. IA fait bien. Cœur = jugement.

**Q : Tâches ingrates ?**
- Pour juniors, oui. Passer des heures à structurer, formuler, vérifier = rigueur.
- Moi, 3 ans, j'ai déjà réflexe. Donc délègue rédaction à IA pour me concentrer fond.
- Vérifie tout. Mais ne perds plus temps sur ce que l'outil fait bien.

**Q : Comment vérifiez ce que produit l'outil ?**
- Trois étapes :
  1. Relecture cohérence : texte tient debout ?
  2. Contrôle croisé sources : infos exactes ?
  3. Question : "oserais défendre en réunion client ?" Si non, je reprends.
- Règle : si je peux pas l'expliquer, je laisse pas.

---

## T3 — Client et valeur

**Q : Confiance client repose sur quoi ?**
- Fiabilité. Cybersécurité, si erreur, conséquences graves.
- Client doit savoir : ils savent ce qu'ils font, prennent pas risques inutiles.
- IA aide à produire vite, pas remplacer jugement expert.
- Confiance = capacité expliquer, justifier, prendre position.

**Q : Diriez-vous au client qu'une partie est IA ?**
- Non. Pas cacher, mais contre-productif. Client paie expertise, pas outil.
- Si dis "c'est l'IA", va se demander pourquoi nous paie.
- Produit final validé par nous => notre responsabilité. Origine du texte secondaire.

**Q : Livrable moins crédible ?**
- Erreur factuelle. Donnée fausse, référence inexacte, concept mal interprété.
- Client perd confiance, et a raison. Cybersécurité, moindre erreur exploitable.
- Vu livrables avec coquilles = collègues pas relu.
- Erreur non vérifiée = faute pro, IA ou main.

**Q : Dans 10 ans, client paiera pour quoi ?**
- Jugement et engagement. IA produira rapports, analyses, scénarios. Pas décision, pas responsabilité.
- Besoin de gens qui disent "voilà ce qu'il faut faire, et pourquoi". Pas voir IA faire ça.

---

## T4 — Gouvernance

**Q : Règles sur usage ?**
- Politique officielle. Copilot approuvé pour usage général, pas données sensibles/classifiées.
- Règle assez générale. Mon contexte = exigences plus fortes. J'applique règle + stricte.
- Prudence.

**Q : Avant livrable client ?**
- Circuit validation. Produis version -> relecteur (manager/expert) -> validation.
- Pas de procédure spécifique pour contenus IA. Consultant s'assure.
- C'est moi qui signe => moi responsable.

**Q : Déjà rencontré contenu faux IA ?**
- Oui, plusieurs fois. Plus récente : synthèse Copilot mal interprété point technique. Rien dramatique, repéré relecture.
- Collègues juniors se font avoir. Envoient livrable, manager voit incohérence. Junior répond "c'est l'IA". Jamais bien fini.

**Q : Règle idéale ?**
- "IA outil d'assistance, pas substitut expertise. Toute production vérifiée et justifiée. En cas de doute, pas utilisée."
- Traçabilité : savoir ce qui a été fait par IA et pas. Pas contrôler, responsabiliser.

---

## T5 — Clôture

**Q : À la place de la direction, première décision ?**
- Formation obligatoire. Pas technique, mais détection erreurs. Risque = consultants ne voient pas les erreurs.
- Processus validation incluant vérification spécifique contenus générés. Comme pour citations, mais systématique.

**Q : Aspect important non abordé ?**
- Dimension juridique. PI, confidentialité, responsabilité. Personne en parle vraiment. Va devenir central.
- Vide juridique à combler.

**Q : Point de vue utile à recueillir ?**
- Responsable juridique cabinet. Comment il perçoit risques : confidentialité, PI, responsabilité.
- Moi vois risques opérationnels, pas légaux. Probablement plus grands qu'on imagine.